\chapter{Common}
\label{chap:common}

This chapter collects definitions and results that are not specific to any one
paper: the multigraph infrastructure, and the general theory of families of
objects, Fulkerson duality, and modulus. Later papers' chapters specialize
this material (e.g.\ to spanning trees) and build their results on top of it.

\section{Multigraphs}

\begin{definition}[Multigraph]
\label{def:multigraph}
\lean{Multigraph}
\leanok
Following the paragraph preceding Section 1.1 of \feupaperref{}: throughout, $G
= (V, E)$ denotes a finite, connected multigraph with vertex set $V$ and edge
set $E$. Parallel edges are allowed; self-loops are not.

To formalize this, a multigraph on vertex type $V$ and edge type $E$ consists of
a map $\mathrm{endpoints} \colon E \to \operatorname{Sym}^2(V)$ assigning to
each edge its (unordered) pair of endpoints. By making the edges a separate type
(rather than a relation on $V$ or a subset of unordered pairs of vertices), we
implicilty allow parallel edges. They are  simply distinct elements of $E$ with
the same image under $\mathrm{endpoints}$. Finiteness, connectedness, and the
absence of self-loops are not built into the type itself. Instead, they are
added explicitly wherever a paper's results need them, matching the
``throughout this paper'' statement.

For weighted multigraphs, $G = (V, E, \sigma)$, we treat these as a multigraph
$G = (V, E)$ together with a weight function $\sigma \in \mathbb R_{>0}^E$.
\end{definition}

\begin{definition}[Spanning tree]
\label{def:spanning-tree}
\uses{def:multigraph}
\lean{Multigraph.IsSpanningTree}
\leanok
A set of edges $T \subseteq E$ is a \emph{spanning tree} of $G$ if it touches
every vertex of $V$ and contains no cycle. In particular, this immediately rules
out parallel edges, which would form a 2-cycle.

This is formalized in \texttt{Multigraph.IsSpanningTree}: $T$ is acyclic (no
loops, no parallel edges within $T$, and the simple graph obtained by forgetting
edge identities is acyclic) and that same simple graph is connected, i.e.\
touches and joins every vertex of $V$.
\end{definition}

\section{Families of objects and Fulkerson duality}

This section follows Section 1.5 of \feupaperref{}. Families of objects generalize
spanning trees (and paths, cuts, etc.) to a common framework, in which the
notion of an admissible density and its Fulkerson dual family are defined.
Spanning trees are recovered as the special case worked out in
Section~\ref{sec:spanning-tree-modulus}.

\begin{definition}[Family of objects]
\label{def:family-of-objects}
\uses{def:multigraph}
Let $G = (V, E, \sigma)$ be a multigraph with edge weights $\sigma \in
\mathbb R_{>0}^E$. A \emph{family of objects} on $G$ is a pair $(\Gamma, \usageN)$
where $\Gamma$ is a countable index set and $\usageN \in \mathbb R_{\ge 0}^{\Gamma
\times E}$ is a \emph{usage matrix}: $\usageN(\gamma, e)$ is the usage of edge $e$ by
object $\gamma$. We identify each $\gamma \in \Gamma$ with its usage vector
$\usageN(\gamma, \cdot)^T \in \mathbb R_{\ge 0}^E$ ($\eta$-space).
\end{definition}

\begin{definition}[Length of an object]
\label{def:length}
\uses{def:family-of-objects}
A point $\rho \in \mathbb R_{\ge 0}^E$ ($\rho$-space) assigns a cost $\rho(e)$
to each edge. The \emph{length} of an object $\gamma \in \Gamma$ with respect
to $\rho$ is
\[
  \ell_\rho(\gamma) := \sum_{e \in E} \usageN(\gamma, e) \rho(e) = \usageN\rho.
\]
\end{definition}

\begin{definition}[Admissible density and Fulkerson dual family]
\label{def:fulkerson-dual}
\uses{def:length}
A density $\rho \in \mathbb R_{\ge 0}^E$ is \emph{admissible} for a family
$(\Gamma, \usageN)$ if $\ell_\rho(\gamma) \ge 1$ for every $\gamma \in \Gamma$,
equivalently $\usageN \rho \ge \mathbf 1$. Write
\[
  \Adm(\Gamma) := \{\rho \in \mathbb R_{\ge 0}^E : \usageN \rho \ge \mathbf 1\}.
\]
The \emph{Fulkerson dual family} of $\Gamma$ is $\widehat\Gamma :=
\operatorname{ext}(\Adm(\Gamma)) \subset \mathbb R_{\ge 0}^E$, the set of
extreme points of $\Adm(\Gamma)$, viewed as a family of objects in
$\rho$-space.
\end{definition}

\begin{theorem}[Fulkerson duality]
\label{thm:fulkerson-duality}
\uses{def:fulkerson-dual}
For any family of objects $\Gamma$, $\widehat{\widehat\Gamma} \subseteq \Gamma$.
\end{theorem}

\section{Modulus of families of objects}

This section follows Section 1.6 of \feupaperref{}.

\begin{definition}[$p$-modulus]
\label{def:modulus}
\uses{def:fulkerson-dual}
Fix $1 \le p < \infty$. The \emph{$p$-modulus} of a family of objects $\Gamma$
on $G = (V, E, \sigma)$ is
\[
  \Mod_{p,\sigma}(\Gamma) := \inf_{\rho \in \Adm(\Gamma)} \sum_{e \in E} \sigma(e) \rho(e)^p.
\]
A density achieving this infimum is called \emph{extremal} (or optimal); it is
unique when $1 < p < \infty$.
\end{definition}

\begin{theorem}[Fulkerson duality for modulus]
\label{thm:modulus-fulkerson-duality}
\uses{def:modulus}
Let $1 < p < \infty$ and let $q$ satisfy $pq = p + q$. Then
\[
  \Mod_{p,\sigma}(\Gamma)^{1/p} \, \Mod_{q,\sigma^{1-q}}(\widehat\Gamma)^{1/q} = 1,
\]
and the extremal density $\eta^*$ for $\Mod_{q,\sigma^{1-q}}(\widehat\Gamma)$ is
related to the extremal density $\rho^*$ for $\Mod_{p,\sigma}(\Gamma)$ by
\[
  \eta^*(e) = \frac{\sigma(e) \rho^*(e)^{p-1}}{\Mod_{p,\sigma}(\Gamma)} \qquad \forall e \in E.
\]
In particular, when $p = q = 2$,
\[
  \Mod_{2,\sigma}(\Gamma) \, \Mod_{2,\sigma^{-1}}(\widehat\Gamma) = 1,
  \qquad
  \eta^*(e) = \frac{\sigma(e)}{\Mod_{2,\sigma}(\Gamma)} \rho^*(e).
\]
\end{theorem}
